%% HOW TO USE THIS TEMPLATE:
 


%请使用xelatex编译
\documentclass[11pt,addpoints,answers]{exam}


% required macros -- get header.tex file from course web page
% !Mode:: "TeX:UTF-8"
% !TEX root=./hw1.tex
\usepackage{xeCJK}    % 写中文要用到
\usepackage{tikz}
\usepackage{amsmath,amsfonts,amssymb,amsthm}
\usepackage{fullpage}
\usepackage{times}
\usepackage{hyperref}
\usepackage{pdfsync}
\usepackage{microtype}
\usepackage{color}
\usepackage{cleveref}
\crefformat{footnote}{#2\footnotemark[#1]#3}
\definecolor{light-gray}{gray}{0.5}

\renewcommand{\tablename}{表}
\renewcommand{\figurename}{图}
\renewcommand{\proof}{证明}

% \theoremstyle{plain}
% \newtheorem{theorem}{定理~}
% \newtheorem{lemma}{引理~}
% \newtheorem{axiom}{公理~}
% \newtheorem{proposition}{命题~}
% \newtheorem{prop}{性质~}
% \newtheorem{corollary}{推论~}
% \newtheorem{conclusion}{结论~}
% \newtheorem{definition}{定义~}
% \newtheorem{conjecture}{猜想~}
% \newtheorem{example}{例~}
% \newtheorem{remark}{注~}
% \newtheorem{algorithm}{算法~}


%%% BLACKBOARD SYMBOLS

% \newcommand{\C}{\ensuremath{\mathbb{C}}}
\newcommand{\D}{\ensuremath{\mathbb{D}}}
\newcommand{\F}{\ensuremath{\mathbb{F}}}
% \newcommand{\G}{\ensuremath{\mathbb{G}}}
\newcommand{\J}{\ensuremath{\mathbb{J}}}
\newcommand{\N}{\ensuremath{\mathbb{N}}}
\newcommand{\Q}{\ensuremath{\mathbb{Q}}}
\newcommand{\R}{\ensuremath{\mathbb{R}}}
\newcommand{\T}{\ensuremath{\mathbb{T}}}
\newcommand{\Z}{\ensuremath{\mathbb{Z}}}
\newcommand{\QR}{\ensuremath{\mathbb{QR}}}

\newcommand{\Zt}{\ensuremath{\Z_t}}
\newcommand{\Zp}{\ensuremath{\Z_p}}
\newcommand{\Zq}{\ensuremath{\Z_q}}
\newcommand{\ZN}{\ensuremath{\Z_N}}
\newcommand{\Zps}{\ensuremath{\Z_p^*}}
\newcommand{\ZNs}{\ensuremath{\Z_N^*}}
\newcommand{\JN}{\ensuremath{\J_N}}
\newcommand{\QRN}{\ensuremath{\QR_{N}}}
\newcommand{\QRp}{\ensuremath{\QR_{p}}}

%%% THEOREM COMMANDS

\theoremstyle{plain}            % following are "theorem" style
\newtheorem{theorem}{定理}[section]
\newtheorem{lemma}[theorem]{引理}
\newtheorem{corollary}[theorem]{推论}
\newtheorem{proposition}[theorem]{命题}
\newtheorem{claim}[theorem]{声明}
\newtheorem{fact}[theorem]{事实}

\theoremstyle{definition}       % following are def style
\newtheorem{definition}[theorem]{定义}
\newtheorem{conjecture}[theorem]{推测}
\newtheorem{example}[theorem]{例}
\newtheorem{protocol}[theorem]{协议}

\theoremstyle{remark}           % following are remark style
\newtheorem{remark}[theorem]{注}
\newtheorem{note}[theorem]{注}
\newtheorem{exercise}[theorem]{练习}

% equation numbering style
\numberwithin{equation}{section}

%%% GENERAL COMPUTING

\newcommand{\bit}{\ensuremath{\set{0,1}}}
\newcommand{\pmone}{\ensuremath{\set{-1,1}}}

% asymptotics
\DeclareMathOperator{\poly}{poly}
\DeclareMathOperator{\polylog}{polylog}
\DeclareMathOperator{\negl}{negl}
\newcommand{\Otil}{\ensuremath{\tilde{O}}}

% probability/distribution stuff
\DeclareMathOperator*{\E}{E}
\DeclareMathOperator*{\Var}{Var}

% sets in calligraphic type
\newcommand{\calA}{\ensuremath{\mathcal{A}}}
\newcommand{\calD}{\ensuremath{\mathcal{D}}}
\newcommand{\calF}{\ensuremath{\mathcal{F}}}
\newcommand{\calH}{\ensuremath{\mathcal{H}}}
\newcommand{\calK}{\ensuremath{\mathcal{K}}}
\newcommand{\calM}{\ensuremath{\mathcal{M}}}
\newcommand{\calX}{\ensuremath{\mathcal{X}}}
\newcommand{\calY}{\ensuremath{\mathcal{Y}}}

% types of indistinguishability
\newcommand{\compind}{\ensuremath{\stackrel{c}{\approx}}}
\newcommand{\statind}{\ensuremath{\stackrel{s}{\approx}}}
\newcommand{\perfind}{\ensuremath{\equiv}}

% font for general-purpose algorithms
\newcommand{\algo}[1]{\ensuremath{\mathsf{#1}}}
% font for general-purpose computational problems
\newcommand{\problem}[1]{\ensuremath{\mathsf{#1}}}
% font for complexity classes
\newcommand{\class}[1]{\ensuremath{\mathsf{#1}}}

% complexity classes and languages
\renewcommand{\P}{\class{P}}
\newcommand{\BPP}{\class{BPP}}
\newcommand{\NP}{\class{NP}}
\newcommand{\coNP}{\class{coNP}}
\newcommand{\AM}{\class{AM}}
\newcommand{\coAM}{\class{coAM}}
\newcommand{\IP}{\class{IP}}

%%% "LEFT-RIGHT" PAIRS OF SYMBOLS

% inner product
\newcommand{\inner}[1]{\langle{#1}\rangle}
\newcommand{\innerfit}[1]{\left\langle{#1}\right\rangle}
% absolute value
\newcommand{\abs}[1]{\lvert{#1}\rvert}
\newcommand{\absfit}[1]{\left\lvert{#1}\right\rvert}
% a set
\newcommand{\set}[1]{\{{#1}\}}
\newcommand{\setfit}[1]{\left\{{#1}\right\}}
% parens
\newcommand{\parens}[1]{({#1})}
\newcommand{\parensfit}[1]{\left({#1}\right)}
% tuple = alias for parens
\newcommand{\tuple}[1]{\parens{#1}}
\newcommand{\tuplefit}[1]{\parensfit{#1}}
% square brackets
\newcommand{\bracks}[1]{[{#1}]}
\newcommand{\bracksfit}[1]{\left[{#1}\right]}
% rounding off
\newcommand{\round}[1]{\lfloor{#1}\rceil}
% floor function
\newcommand{\floor}[1]{\lfloor{#1}\rfloor}
% ceiling function
\newcommand{\ceil}[1]{\lceil{#1}\rceil}
% length of a string
\newcommand{\len}[1]{\lvert{#1}\rvert}
\newcommand{\lenfit}[1]{\left\lvert{#1}\right\rvert}
% length of some vector, element
\newcommand{\length}[1]{\lVert{#1}\rVert}
\newcommand{\lengthfit}[1]{\left\lVert{#1}\right\rVert}

%%% CRYPTO-RELATED NOTATION

% KEYS AND RELATED

\newcommand{\key}[1]{\ensuremath{#1}}

\newcommand{\pk}{\key{pk}}
\newcommand{\vk}{\key{vk}}
\newcommand{\sk}{\key{sk}}
\newcommand{\mpk}{\key{mpk}}
\newcommand{\msk}{\key{msk}}
\newcommand{\fk}{\key{fk}}
\newcommand{\id}{id}
\newcommand{\keyspace}{\ensuremath{\mathcal{K}}}
\newcommand{\msgspace}{\ensuremath{\mathcal{M}}}
\newcommand{\ctspace}{\ensuremath{\mathcal{C}}}
\newcommand{\tagspace}{\ensuremath{\mathcal{T}}}
\newcommand{\idspace}{\ensuremath{\mathcal{ID}}}

\newcommand{\concat}{\ensuremath{\|}}

% GAMES

% advantage
\newcommand{\advan}{\ensuremath{\mathbf{Adv}}}

% different attack models
\newcommand{\attack}[1]{\ensuremath{\text{#1}}}

\newcommand{\atk}{\attack{atk}} % dummy attack
\newcommand{\indcpa}{\attack{ind-cpa}}
\newcommand{\indcca}{\attack{ind-cca}}
\newcommand{\anocpa}{\attack{ano-cpa}} % anonymous
\newcommand{\anocca}{\attack{ano-cca}}
\newcommand{\euacma}{\attack{eu-acma}} % forgery: adaptive chosen-message
\newcommand{\euscma}{\attack{eu-scma}} % forgery: static chosen-message
\newcommand{\suacma}{\attack{su-acma}} % strongly unforgeable

% ADVERSARIES
\newcommand{\attacker}[1]{\ensuremath{\mathcal{#1}}}

\newcommand{\Adv}{\attacker{A}}
\newcommand{\AdvA}{\attacker{A}}
\newcommand{\AdvB}{\attacker{B}}
\newcommand{\Dist}{\attacker{D}}
\newcommand{\Sim}{\attacker{S}}
\newcommand{\Ora}{\attacker{O}}
\newcommand{\Inv}{\attacker{I}}
\newcommand{\For}{\attacker{F}}

% CRYPTO SCHEMES

\newcommand{\scheme}[1]{\ensuremath{\text{#1}}}

% pseudorandom stuff
\newcommand{\prg}{\algo{PRG}}
\newcommand{\prf}{\algo{PRF}}
\newcommand{\prp}{\algo{PRP}}

% symmetric-key cryptosystem
\newcommand{\skc}{\scheme{SKC}}
\newcommand{\skcgen}{\algo{Gen}}
\newcommand{\skcenc}{\algo{Enc}}
\newcommand{\skcdec}{\algo{Dec}}

% public-key cryptosystem
\newcommand{\pkc}{\scheme{PKC}}
\newcommand{\pkcgen}{\algo{Gen}}
\newcommand{\pkcenc}{\algo{Enc}} % can also use \kemenc and \kemdec
\newcommand{\pkcdec}{\algo{Dec}}

% digital signatures
\newcommand{\sig}{\scheme{SIG}}
\newcommand{\siggen}{\algo{Gen}}
\newcommand{\sigsign}{\algo{Sign}}
\newcommand{\sigver}{\algo{Ver}}

% message authentication code
\newcommand{\mac}{\scheme{MAC}}
\newcommand{\macgen}{\algo{Gen}}
\newcommand{\mactag}{\algo{Tag}}
\newcommand{\macver}{\algo{Ver}}

% key-encapsulation mechanism
\newcommand{\kem}{\scheme{KEM}}
\newcommand{\kemgen}{\algo{Gen}}
\newcommand{\kemenc}{\algo{Encaps}}
\newcommand{\kemdec}{\algo{Decaps}}

% identity-based encryption
\newcommand{\ibe}{\scheme{IBE}}
\newcommand{\ibesetup}{\algo{Setup}}
\newcommand{\ibeext}{\algo{Ext}}
\newcommand{\ibeenc}{\algo{Enc}}
\newcommand{\ibedec}{\algo{Dec}}

% hierarchical IBE (as key encapsulation)
\newcommand{\hibe}{\scheme{HIBE}}
\newcommand{\hibesetup}{\algo{Setup}}
\newcommand{\hibeext}{\algo{Extract}}
\newcommand{\hibeenc}{\algo{Encaps}}
\newcommand{\hibedec}{\algo{Decaps}}

% binary tree encryption (as key encapsulation)
\newcommand{\bte}{\scheme{BTE}}
\newcommand{\btesetup}{\algo{Setup}}
\newcommand{\bteext}{\algo{Extract}}
\newcommand{\bteenc}{\algo{Encaps}}
\newcommand{\btedec}{\algo{Decaps}}

% trapdoor functions
\newcommand{\tdf}{\scheme{TDF}}
\newcommand{\tdfgen}{\algo{Gen}}
\newcommand{\tdfeval}{\algo{Eval}}
\newcommand{\tdfinv}{\algo{Invert}}
\newcommand{\tdfver}{\algo{Ver}}

%%% PROTOCOLS

\newcommand{\out}{\text{out}}
\newcommand{\view}{\text{view}}

%%% COMMANDS FOR LECTURES/HOMEWORKS

\newcommand{\lecheader}{%
  \chead{\large \textbf{Lecture \lecturenum\\\lecturetopic}}

  \lhead{\small
    \textbf{\href{http://bhpan.buaa.edu.cn}{研究生课程：计算复杂性}\\2023年秋季}}

  \rhead{\small \textbf{主讲: 张宗洋
      }

  \setlength{\headheight}{20pt}
  \setlength{\headsep}{16pt}
}}

\newcommand{\hwheader}{%
  \chead{\Large \textbf{作业 \hwnum}}

  \lhead{\small
    \textbf{{计算复杂性}\\2025年秋季}}

  \rhead{\small \textbf{主讲: 张宗洋
      \\姓名: \studentname}}

  \setlength{\headheight}{20pt}
  \setlength{\headsep}{16pt}
  
  \headrule
}

\newcommand{\examheader}{%
  \chead{\Large \textbf{测试 \\ \duedate}}
  
\lhead{\small
    \textbf{\href{http://bhpan.buaa.edu.cn}{计算复杂性}\\2024年秋季}}

  \rhead{\small \textbf{主讲: 张宗洋
      \\姓名: \studentname}}

  \setlength{\headheight}{20pt}
  \setlength{\headsep}{16pt}
  
  \headrule
}
  
	
\newcommand{\hint}[1]{\href{http://www.cims.nyu.edu/~regev/cgi-bin/hints/index.html?#1}{{\textcolor{light-gray}{I need a hint! (ID #1)}}}}

\newcommand{\hinttext}[2]{\href{http://www.cims.nyu.edu/~regev/cgi-bin/hints/index.html?#1}{{\textcolor{light-gray}{#2 (ID #1)}}}}

\newcommand{\hintcost}[2]{\href{http://www.cims.nyu.edu/~regev/cgi-bin/hints/index.html?#1}{{\textcolor{light-gray}{I need a hint for #2 points! (ID #1)}}}}
	
\newcommand{\moreinfo}[1]{\href{http://www.cims.nyu.edu/~regev/cgi-bin/hints/index.html?#1}{{\textcolor{light-gray}{I'm done solving and want to know more! (ID #1)}}}}

% VARIABLES

\newcommand{\hwnum}{5}
\newcommand{\duedate}{12月17日晚23:59前} % changing this does not change the
                                % actual due date :)
\newcommand{\studentname}{}

% END OF SUPPLIED VARIABLES

\hwheader                       % execute homework commands

\begin{document}

\pagestyle{head}                % put header on  

\pagestyle{head}                % put header on  

\begin{itemize}
	\item 请将PDF文件发送至邮箱：complexity\_buaa@163.com
	\item 文件命名规则是``学号-姓名-第五次作业.pdf"
	\item 务必独立自主解题
	\item 作业截止时间为 \textbf{\duedate}.
\end{itemize}

\medskip


 

% QUESTIONS START HERE.  PROVIDE SOLUTIONS WITHIN THE "solution"
% ENVIRONMENTS FOLLOWING EACH QUESTION.

\begin{questions}


\question P151, 5.1题。
证明$EQ_{\textit{CFG}}$是不可判定的。
\begin{solution}
    通过归约来证明。已知：
    \[
    ALL_{\text{CFG}} = \{ \langle G \rangle \mid G \text{是CFG，且} L(G) = \Sigma^* \}
    \]
    是不可判定的。
    
    现在构造从$ALL_{\text{CFG}}$到$EQ_{\text{CFG}}$的归约。给定任意CFG $G$，构造两个CFG $G_1$和$G_2$如下：
    \begin{enumerate}
        \item 令$G_1 = G$。
        \item 构造$G_2$使得$L(G_2) = \Sigma^*$。这总是可行的，因为$\Sigma^*$是正则语言，因此也是上下文无关语言。
    \end{enumerate}
    
    定义归约函数$f$为：
    \[
    f(\langle G \rangle) = \langle G_1, G_2 \rangle = \langle G, G_{\Sigma^*} \rangle
    \]
    其中$G_{\Sigma^*}$是生成$\Sigma^*$的固定CFG。
    
    下面证明归约的正确性：
    \begin{enumerate}
        \item 如果$\langle G \rangle \in ALL_{\text{CFG}}$，则$L(G) = \Sigma^* = L(G_{\Sigma^*})$，因此$\langle G, G_{\Sigma^*} \rangle \in EQ_{\text{CFG}}$。
        \item 如果$\langle G \rangle \notin ALL_{\text{CFG}}$，则$L(G) \neq \Sigma^*$，而$L(G_{\Sigma^*}) = \Sigma^*$，因此$L(G) \neq L(G_{\Sigma^*})$，所以$\langle G, G_{\Sigma^*} \rangle \notin EQ_{\text{CFG}}$。
    \end{enumerate}
    
    因此，$ALL_{\text{CFG}} \leq_m EQ_{\text{CFG}}$。由于$ALL_{\text{CFG}}$是不可判定的，所以$EQ_{\text{CFG}}$也是不可判定的。

\end{solution}



\question P151, 5.2题。
证明$EQ_{\textit{CFG}}$是补图灵可识别的。
\begin{solution}
	设$G_1$和$G_2$是两个上下文无关文法（CFG），其字母表为$\Sigma$。
    我们需要证明存在一个图灵机可以识别：
    \[
    \overline{EQ_{\text{CFG}}} = \{\langle G_1, G_2 \rangle \mid L(G_1) \neq L(G_2)\}
    \]
    \noindent \textbf{构造识别机$M$：}
    \begin{enumerate}[label=步骤\arabic*:, leftmargin=*]
        \item 输入为$\langle G_1, G_2 \rangle$。
        \item 枚举所有字符串$x \in \Sigma^*$（按标准顺序，如短lex序）。
        \item 对每个枚举的字符串$x$，并行地进行以下测试：
        \begin{itemize}
            \item 测试$x \in L(G_1)$且$x \notin L(G_2)$
            \item 测试$x \in L(G_2)$且$x \notin L(G_1)$
        \end{itemize}
        \item 对于每个$x$，利用两个CFG的成员资格判定算法（由于CFG的成员资格问题是可判定的，存在算法判断$x \in L(G)$）。
        \item 如果发现某个$x$使得上述任一条件成立，则$M$接受$\langle G_1, G_2 \rangle$。
    \end{enumerate}
    
    \noindent \textbf{正确性分析：}
    \begin{itemize}
        \item 如果$L(G_1) \neq L(G_2)$，则必然存在某个$x$在对称差中，因此$M$最终会发现这样的$x$并接受。
        \item 如果$L(G_1) = L(G_2)$，则对所有$x$，两个测试都会失败，$M$将永不停机（即永远搜索下去）。这符合识别机的定义：对于不在语言中的输入，$M$可能永不停机。
    \end{itemize}
    
    因此，$M$识别了$\overline{EQ_{\text{CFG}}}$，即$\overline{EQ_{\text{CFG}}}$是图灵可识别的。
\end{solution}







\question P151,5.4题。
如果$A\leq _{m} B$且$B$是一个正则语言，这是否蕴涵着$A$也是一个正则语言？为什么？
\begin{solution}
    	题目中 $A \leq_m B$ 仅要求 $f$ 是可计算函数，不要求 $f$ 是正则的或线性的。因此即使 $B$ 是正则语言，$A$ 也可以是任意图灵可识别的语言（甚至可判定语言），包括非正则语言。例如：
    \begin{itemize}
        \item $B = \emptyset$（正则）
        \item $A$ 是任意语言，定义 $f(x) = \text{某个不在 } B \text{ 中的字符串若 } x \notin A$，\\否则 $f(x) = \text{某个在 } B \text{ 中的字符串}$。
    \end{itemize}
    实际上，如果 $B$ 是正则语言且 $B \neq \emptyset$ 且 $B \neq \Sigma^*$，那么仍可构造非正则的 $A$ 通过一个复杂的可计算 $f$ 归约到 $B$。例如，令 $A$ 为不可判定的语言，$B = \{0\}$（正则），则存在可计算函数 $f$ 将 $A$ 归约到 $B$，只需将 $A$ 的成员映射到 $0$，非成员映射到 $1$（但 $1 \notin B$，所以 $x \notin A \Rightarrow f(x) \notin B$）。这可行，因为 $B$ 是正则的。
    
    结论：$A \leq_m B$ 且 $B$ 是正则语言 \textbf{不} 蕴涵 $A$ 是正则语言，因为 many-one 归约只要求 $f$ 可计算，不保留正则性。

\end{solution}



\question P151, 5.5题。
证明$A_{\textit{TM}}$无法映射归约到$E_{\textit{TM}}$。换句话说，也就是证明没有可计算函数可以将$A_{\textit{TM}}$归约为$E_{\textit{TM}}$。（提示：用$A_{\textit{TM}}$和$E_{\textit{TM}}$的一些已知的矛盾和事实来证明。）
\begin{solution}
	已知：$E_{\text{TM}}$ 不是图灵可识别的。
    % \begin{enumerate}
    %     \item $E_{\text{TM}}$ 不是图灵可识别的。
    %     % \item $\overline{E_{\text{TM}}}$ 是图灵可识别的。
    % \end{enumerate}
    
    如果 $A_{\text{TM}} \leq_m E_{\text{TM}}$，那么因为 $E_{\text{TM}}$ 不是可识别的，所以 $A_{\text{TM}}$ 也必须是不可识别的（因为可识别性在映射归约下反向传递：如果 $A$ 可识别且 $A \leq_m B$，则 $B$ 可识别；等价地，如果 $B$ 不可识别且 $A \leq_m B$，则 $A$ 不可识别）。
    
    但已知 $A_{\text{TM}}$ 是可识别的，矛盾。
    
    因此，假设错误，不存在这样的可计算函数 $f$。
    故 $A_{\text{TM}} \not\leq_m E_{\text{TM}}$
\end{solution}


\question P151,5.7题。
证明：如果$A$是图灵可识别的，且$A\leq _{m} \overline{A}$，则A是可判定的。
\begin{solution}
	设 $f$ 是归约函数，即存在可计算函数 $f$ 使得
    \[
    x \in A \iff f(x) \in \overline{A} \quad \text{对任意输入 } x
    \]
    等价地，
    \[
    x \in A \iff f(x) \notin A
    \]
    
    由于 $A$ 是图灵可识别的，存在图灵机 $M$ 识别 $A$，即 $L(M) = A$。  
    同样地，$\overline{A}$ 也是图灵可识别的吗？不一定，但我们可以利用归约的性质。
    
    考虑以下判定 $A$ 的算法：
    
    \noindent \textbf{算法构造：}
    输入 $x$：
    \begin{enumerate}
        \item 计算 $f(x)$。
        \item 并行运行以下两个过程：
        \begin{itemize}
            \item 过程1：在输入 $x$ 上运行 $A$ 的识别器 $M$（如果 $x \in A$，此过程会接受）。
            \item 过程2：在输入 $f(x)$ 上运行 $A$ 的识别器 $M$（如果 $f(x) \in A$，此过程会接受）。
        \end{itemize}
        \item 等待其中一个过程接受：
        \begin{itemize}
            \item 如果过程1接受，则 $x \in A$，输出“接受”。
            \item 如果过程2接受，则 $f(x) \in A$，根据归约性质 $x \notin A$，输出“拒绝”。
        \end{itemize}
    \end{enumerate}
    
    \noindent \textbf{正确性证明：}
    我们需要证明上述算法一定会终止，并且正确判定 $x$ 是否属于 $A$。
    
    根据归约性质 $x \in A \iff f(x) \notin A$，所以 $x$ 和 $f(x)$ 恰好有一个属于 $A$。因此：
    \begin{itemize}
        \item 如果 $x \in A$，则过程1会接受（因为 $M$ 识别 $A$），过程2不会接受（因为 $f(x) \notin A$）。
        \item 如果 $x \notin A$，则 $f(x) \in A$，过程2会接受，过程1不会接受。
    \end{itemize}
    
    所以无论如何，两个过程中总有一个会接受，算法必然在有限步内停止（因为 $M$ 是识别器，对属于 $A$ 的输入一定会接受）。当算法停止时，它能正确输出 $x \in A$ 或 $x \notin A$。
    
    因此， $A$ 是可判定的。
\end{solution}



\question P151,5.10题。
证明当且仅当$A\leq _{m} A_\text{TM}$时，A是图灵可识别的。
\begin{solution}
	\noindent \textbf{（充分性 $\Rightarrow$）} 假设 $A$ 是图灵可识别的。  
    存在图灵机 $M_A$ 识别 $A$，即 $L(M_A) = A$。定义归约函数 $f$ 如下：
    \[
    f(x) = \langle M_A, x \rangle
    \]
    这里 $f$ 是可计算的，因为给定 $x$，我们可以输出描述 $M_A$ 的编码和 $x$ 的字符串 $\langle M_A, x \rangle$。
    
    现在验证归约的正确性：
    \[
    x \in A \iff M_A \text{ 接受 } x \iff \langle M_A, x \rangle \in A_{\text{TM}}
    \]
    因此 $A \leq_m A_{\text{TM}}$。
    
    \vspace{0.5em}
    \noindent \textbf{（必要性 $\Leftarrow$）} 假设 $A \leq_m A_{\text{TM}}$。  
    由于 $A_{\text{TM}}$ 是图灵可识别的（存在通用图灵机识别它），且图灵可识别性在映射归约下保持：若 $A \leq_m B$ 且 $B$ 可识别，则 $A$ 可识别。  
    因此 $A$ 是图灵可识别的。
\end{solution}


\question P151,5.11题。
证明当且仅当$A\leq _{m} 0*1*$时，A是可判定的。
\begin{solution}

    \noindent \textbf{1. ($\Rightarrow$)：若 $A$ 可判定，则 $A \leq_m 0^*1^*$。}
    
    设 $A$ 是可判定的，则存在图灵机 $M$ 判定 $A$。构造一个可计算函数 $f$，使得
    \[
    x \in A \iff f(x) \in 0^*1^*.
    \]
    定义 $f$ 如下：
    \[
    f(x) = 
    \begin{cases}
    0 & \text{如果 } x \in A, \\
    1 & \text{如果 } x \notin A.
    \end{cases}
    \]
    由于 $M$ 判定 $A$，我们可以计算 $f(x)$：先运行 $M$ 在 $x$ 上，根据 $M$ 接受或拒绝，分别输出 $0$ 或 $1$。
    
    现在验证归约的正确性：
    \begin{itemize}
        \item 如果 $x \in A$，则 $f(x) = 0 \in 0^*1^*$。
        \item 如果 $x \notin A$，则 $f(x) = 1$。注意 $1 \in 0^*1^*$（取 $i=0, j=1$），因此 $f(x) \in 0^*1^*$ 仍然成立。
    \end{itemize}
    这意味着对所有 $x$，$f(x) \in 0^*1^*$，即 $f$ 将所有输入映射到 $0^*1^*$ 中。这导致归约无法区分 $A$ 的成员与非成员。事实上，上述构造有误，因为它使得 $x \in A$ 和 $x \notin A$ 都映射到 $0^*1^*$ 中，不符合归约必须保持成员与非成员区分的要求。
    
    因此需要修改构造，使得 $f(x) \in 0^*1^*$ 当且仅当 $x \in A$。一种正确的构造如下：
    
    取 $y_0 \in 0^*1^*$（例如 $y_0 = \varepsilon$），$y_1 \notin 0^*1^*$（例如 $y_1 = 10$，因为 $10 \notin 0^*1^*$）。定义：
    \[
    f(x) = 
    \begin{cases}
    y_0 & \text{如果 } x \in A, \\
    y_1 & \text{如果 } x \notin A.
    \end{cases}
    \]
    由于 $A$ 可判定，$f$ 可计算，并且
    \[
    x \in A \iff f(x) = y_0 \in 0^*1^*,
    \]
    \[
    x \notin A \iff f(x) = y_1 \notin 0^*1^*.
    \]
    因此 $A \leq_m 0^*1^*$。
    
    \vspace{1em}
    \noindent \textbf{2. ($\Leftarrow$)：若 $A \leq_m 0^*1^*$，则 $A$ 可判定。}
    
    假设存在可计算函数 $f$ 使得 $x \in A \iff f(x) \in 0^*1^*$。由于 $0^*1^*$ 是正则语言，因此是可判定的（存在算法判定一个字符串是否属于 $0^*1^*$）。要判定 $A$，可对输入 $x$ 执行如下算法：
    \begin{enumerate}
        \item 计算 $f(x)$。
        \item 判定 $f(x) \in 0^*1^*$ 是否成立（这是可判定的）。
        \item 若成立则接受 $x$，否则拒绝 $x$。
    \end{enumerate}
    该算法总是终止，因为 $f$ 可计算且 $0^*1^*$ 可判定。根据归约性质，它正确判定 $A$。因此 $A$ 是可判定的。
\end{solution}


\question P151, 5.12题。
设$J=\left\{w|\text{对于某个}x\in A_{\textit{TM}}\text{有}w=0x\text{，或对于某个}y\in \overline{A_{\textit{TM}}} \text{有} w=1y \right\}$。证明$J$和$\overline J$都不是图灵可识别的。
\begin{solution}

    证明更强的结论：如果 $J$ 可识别，则 $\overline{A_{\text{TM}}}$ 可识别；如果 $\overline{J}$ 可识别，则 $\overline{A_{\text{TM}}}$ 可识别。但已知 $A_{\text{TM}}$ 可识别，若 $\overline{A_{\text{TM}}}$ 也可识别，则 $A_{\text{TM}}$ 可判定，矛盾。因此 $J$ 与 $\overline{J}$ 均不可识别。
    
    具体地：
    \begin{itemize}
    \item 若 $J$ 可识别，则通过映射 $f(\langle M,w\rangle) = 1\langle M,w\rangle$，可得 $\overline{A_{\text{TM}}} \leq_m J$，从而 $\overline{A_{\text{TM}}}$ 可识别。
    \item 若 $\overline{J}$ 可识别，则通过映射 $g(\langle M,w\rangle) = 0\langle M,w\rangle$，可得 $\overline{A_{\text{TM}}} \leq_m \overline{J}$，从而 $\overline{A_{\text{TM}}}$ 可识别。
    \end{itemize}
    但 $A_{\text{TM}}$ 可识别，$\overline{A_{\text{TM}}}$ 不可识别（否则 $A_{\text{TM}}$ 可判定）。矛盾。因此 $J$ 与 $\overline{J}$ 均不可识别。
\end{solution}


\question P151,5.13题。
给出一个不可判定语言B的例子，使得$B\leq _{m} \overline{B}$。
\begin{solution}
	一个标准例子是：
    \[
    B = A_{\mathrm{TM}} \oplus \overline{A_{\mathrm{TM}}} = \{0x \mid x \in A_{\mathrm{TM}}\} \cup \{1x \mid x \in \overline{A_{\mathrm{TM}}}\}.
    \]
    这正是5.12题中的 $J$。已经证明 $J$ 和 $\overline{J}$ 都不可识别，因此 $J$ 不可判定。而且显然有 $J \leq_m \overline{J}$，只需定义 $f(0x) = 1x$，$f(1x) = 0x$ 即可，这满足：
    \[
    w \in J \iff f(w) \in \overline{J}.
    \]
    函数 $f$ 可计算，因此 $J$ 即为所求的不可判定语言 $B$。
\end{solution}



 




\question P152, 5.16题。
\textbf{赖斯定理} \quad 设$P$是任何图灵机的语言的非平凡属性。证明确定某一图灵机的语言是否具有属性$P$这一问题是不可判定的。更正式地说，设$P$是一个语言，它由图灵机的描述组成，且$P$满足下列两个条件：首先，$P$是非平凡的——它包含某些图灵机的描述，但不是所有的；其次，$P$是图灵机的语言的属性——无论何时$L(M_{1})=L(M_{2})$，当且仅当$\left \langle M_{2} \right \rangle \in P $时，有$\left \langle M_{1} \right \rangle \in P $，此处的$M_{1} $和$M_{2} $是任意图灵机。证明$P$是一个不可判定语言。
\begin{solution}

	\subsection*{证明思路}
    通过归约 $A_{\mathrm{TM}} \leq_m P$ 或 $A_{\mathrm{TM}} \leq_m \overline{P}$ 来证明。由于 $A_{\mathrm{TM}}$ 不可判定，因此 $P$ 不可判定。需要根据 $P$ 是否包含空语言 $\emptyset$ 分两种情况构造归约。
    
    \subsection*{证明细节}
    设 $M_P$ 是一台语言属于 $P$ 的图灵机（由非平凡性，这样的 $M_P$ 存在），设 $M_{\emptyset}$ 是一台语言为空集（即 $L(M_{\emptyset}) = \emptyset$）且 $\langle M_{\emptyset} \rangle \notin P$ 的图灵机（由非平凡性，若空语言在 $P$ 中，则取另一不在 $P$ 中的语言对应的机器）。下面分两种情况。
    
    \noindent \textbf{情况1：空语言 $\emptyset \notin P$（即 $\langle M_{\emptyset} \rangle \notin P$）。} \\
    构造从 $A_{\mathrm{TM}}$ 到 $P$ 的归约 $f$。
    
    对输入 $\langle M, w \rangle$，定义 $f(\langle M, w \rangle) = \langle M' \rangle$，其中 $M'$ 为如下图灵机：

    \[
    M' = \text{“对输入 } x:
    ]\
    \begin{enumerate}
        \item 模拟 $M$ 在输入 $w$ 上运行。
        \item 如果 $M$ 接受 $w$，则模拟 $M_P$ 在输入 $x$ 上运行并接受当且仅当 $M_P$ 接受。
        \item 如果 $M$ 拒绝 $w$，则拒绝 $x$（或循环）。”
    \end{enumerate}

    更准确地说，步骤2中若 $M$ 接受 $w$，则 $M'$ 完全按照 $M_P$ 的行为运行；步骤3中若 $M$ 拒绝 $w$，则 $M'$ 直接拒绝所有 $x$（或进入循环，使 $L(M') = \emptyset$）。
    
    \noindent \textbf{分析 $M'$ 的语言：}
    \begin{itemize}
        \item 若 $M$ 接受 $w$，则 $M'$ 的行为与 $M_P$ 相同，因此 $L(M') = L(M_P)$，从而 $\langle M' \rangle \in P$（因为 $P$ 是语言属性）。
        \item 若 $M$ 不接受 $w$（拒绝或不停机），则 $M'$ 不会执行到步骤2，对所有 $x$ 都拒绝（或不停机），因此 $L(M') = \emptyset$，从而 $\langle M' \rangle \notin P$（因为 $\emptyset \notin P$）。
    \end{itemize}
    因此，
    \[
    \langle M, w \rangle \in A_{\mathrm{TM}} \iff L(M') = L(M_P) \iff \langle M' \rangle \in P.
    \]
    故 $f$ 实现了归约 $A_{\mathrm{TM}} \leq_m P$，从而 $P$ 不可判定。
    
    \vspace{1em}
    \noindent \textbf{情况2：空语言 $\emptyset \in P$（即 $\langle M_{\emptyset} \rangle \in P$）。} \\
    此时考虑 $\overline{P}$，显然 $\emptyset \notin \overline{P}$。由情况1的结论（将 $P$ 换为 $\overline{P}$）可知 $\overline{P}$ 不可判定，从而 $P$ 也不可判定。
    
    更直接地，我们构造从 $A_{\mathrm{TM}}$ 到 $\overline{P}$ 的归约。取 $M_{\text{nonP}}$ 为一台语言不在 $P$ 中的图灵机（由非平凡性存在）。对输入 $\langle M, w \rangle$，构造 $M''$：
    \[
    M'' = \text{“对输入 } x:
     \]
    \begin{enumerate}
        \item 模拟 $M$ 在 $w$ 上运行。
        \item 如果 $M$ 接受 $w$，则模拟 $M_{\text{nonP}}$ 在 $x$ 上运行并接受当且仅当 $M_{\text{nonP}}$ 接受。
        \item 如果 $M$ 拒绝 $w$，则拒绝 $x$。”
    \end{enumerate}

    分析：
    \begin{itemize}
        \item 若 $M$ 接受 $w$，则 $L(M'') = L(M_{\text{nonP}})$，故 $\langle M'' \rangle \notin P$，即 $\langle M'' \rangle \in \overline{P}$。
        \item 若 $M$ 不接受 $w$，则 $L(M'') = \emptyset$，由于 $\emptyset \in P$，故 $\langle M'' \rangle \in P$，即 $\langle M'' \rangle \notin \overline{P}$。
    \end{itemize}
    因此 $\langle M, w \rangle \in A_{\mathrm{TM}} \iff \langle M'' \rangle \in \overline{P}$，即 $A_{\mathrm{TM}} \leq_m \overline{P}$，故 $\overline{P}$ 不可判定，从而 $P$ 也不可判定。
    
    \subsection*{总结}
    无论空语言是否属于 $P$，我们都可以构造从 $A_{\mathrm{TM}}$ 到 $P$ 或 $\overline{P}$ 的归约，从而证明 $P$ 是不可判定的。这就完成了赖斯定理的证明。
    
\end{solution}


\question P152, 5.18题。
使用问题5.16中的赖斯定理，证明下列语言是不可判定的：
\begin{parts} 
	\part   $INFINITE_{TM}=\{\langle M\rangle|M\text{是一个图灵机，且}L(M)\text{是一个无限语言}\}$.
	\begin{solution}
		\begin{itemize}
    \item \textbf{非平凡性：} 
    \begin{enumerate}
        \item 存在图灵机 $M_1$ 使得 $L(M_1)$ 是无限的（例如 $M_1$ 接受所有串，则 $L(M_1) = \Sigma^*$ 无限），故 $\langle M_1 \rangle \in INFINITE_{\mathrm{TM}}$。
        \item 存在图灵机 $M_2$ 使得 $L(M_2)$ 是有限的（例如 $M_2$ 拒绝所有串，则 $L(M_2) = \emptyset$ 有限），故 $\langle M_2 \rangle \notin INFINITE_{\mathrm{TM}}$。
    \end{enumerate}
    \item \textbf{语言属性：} 若 $L(M) = L(M')$，则 $L(M)$ 无限当且仅当 $L(M')$ 无限，因此 $\langle M \rangle \in INFINITE_{\mathrm{TM}}$ 当且仅当 $\langle M' \rangle \in INFINITE_{\mathrm{TM}}$。
\end{itemize}
由赖斯定理，$INFINITE_{\mathrm{TM}}$ 不可判定。
	\end{solution}
	
	\part 证明 $\{\langle M\rangle|M\text{是一个图灵机,且}1011\in L(M)\}$
	\begin{solution}
		记该语言为 $P_{1011}$。
\begin{itemize}
    \item \textbf{非平凡性：} 
    \begin{enumerate}
        \item 存在图灵机 $M_1$ 接受 $1011$（例如设计 $M_1$ 仅接受 $1011$），则 $\langle M_1 \rangle \in P_{1011}$。
        \item 存在图灵机 $M_2$ 不接受 $1011$（例如 $M_2$ 拒绝所有串），则 $\langle M_2 \rangle \notin P_{1011}$。
    \end{enumerate}
    \item \textbf{语言属性：} 若 $L(M) = L(M')$，则 $1011 \in L(M)$ 当且仅当 $1011 \in L(M')$，因此 $\langle M \rangle \in P_{1011}$ 当且仅当 $\langle M' \rangle \in P_{1011}$。
\end{itemize}
由赖斯定理，$P_{1011}$ 不可判定。
	\end{solution}
	
	\part 证明 $ALL_\text{TM}=\{\langle M\rangle|M\text{是一个图灵机,且}L(M)=\Sigma^*\}$
	\begin{solution}
		\begin{itemize}
    \item \textbf{非平凡性：} 
    \begin{enumerate}
        \item 存在图灵机 $M_1$ 使得 $L(M_1) = \Sigma^*$（例如接受所有输入的图灵机），故 $\langle M_1 \rangle \in ALL_{\mathrm{TM}}$。
        \item 存在图灵机 $M_2$ 使得 $L(M_2) \neq \Sigma^*$（例如拒绝所有输入的图灵机），故 $\langle M_2 \rangle \notin ALL_{\mathrm{TM}}$。
    \end{enumerate}
    \item \textbf{语言属性：} 若 $L(M) = L(M')$，则 $L(M) = \Sigma^*$ 当且仅当 $L(M') = \Sigma^*$，因此 $\langle M \rangle \in ALL_{\mathrm{TM}}$ 当且仅当 $\langle M' \rangle \in ALL_{\mathrm{TM}}$。
\end{itemize}
由赖斯定理，$ALL_{\mathrm{TM}}$ 不可判定。
	\end{solution}
	
\end{parts}







\question P152, 5.25题。

设$T=\{\langle M\rangle|M\text{是一个图灵机,每当}M\text{接受}\omega \text{时，}M\text{也接受}\omega^{R}\}$，证明T是不可判定的。

\begin{solution}
	应用赖斯定理，因为 $T$ 是图灵机语言的非平凡性质，故 $T$ 不可判定。
\end{solution}




\question P152, 5.26题。
考虑这样的问题：一个双带图灵机，当它在输入$w$上运行时，检查它是否在第二条带子上写下一个非空白符。将这个问题形式化为一个语言，并证明它是不可判定的。
\begin{solution}
	令 $M'$ 为如下双带图灵机：
\begin{itemize}
    \item 第一条带：初始内容为 $w$，用于模拟 $M$ 在 $w$ 上的运行（作为 $M$ 的单一工作带）。
    \item 第二条带：初始全为空白，仅用于在 $M$ 接受时写入标记。
    \item $M'$ 的转移函数确保：
        \begin{enumerate}
            \item 在模拟过程中，$M'$ 不在第二条带上写入任何符号（即保持读写头不动，或只读不写）。
            \item 如果模拟中 $M$ 进入接受状态，则 $M'$ 在第二条带的当前位置写入一个非空白符号 $X$，然后停机。
            \item 如果模拟中 $M$ 进入拒绝状态，则 $M'$ 直接停机，且不触碰第二条带（保持全空白）。
            \item 如果 $M$ 永不停机，则 $M'$ 也永不停机，且第二条带保持全空白。
        \end{enumerate}
\end{itemize}

\subsection*{归约函数}
定义归约函数 $f(\langle M, w \rangle) = \langle M', w \rangle$，其中 $M'$ 是如上构造的双带图灵机（其描述可从 $\langle M, w \rangle$ 可计算地构造）。

\subsection*{正确性证明}
\begin{itemize}
    \item 若 $\langle M, w \rangle \in A_{\text{TM}}$，即 $M$ 接受 $w$，则 $M'$ 在模拟过程中会进入 $M$ 的接受状态，从而在第二条带上写入 $X$，因此 $\langle M', w \rangle \in L_{\text{write}}$。
    \item 若 $\langle M, w \rangle \notin A_{\text{TM}}$，则有两种情况：
        \begin{enumerate}
            \item $M$ 拒绝 $w$：则 $M'$ 模拟 $M$ 直到拒绝状态，然后停机，且从未在第二条带上写入非空白符号，故 $\langle M', w \rangle \notin L_{\text{write}}$。
            \item $M$ 在 $w$ 上永不停机：则 $M'$ 也永不停机，且永远不会在第二条带上写入非空白符号，故 $\langle M', w \rangle \notin L_{\text{write}}$。
        \end{enumerate}
    \end{itemize}
    因此，
    \[
    \langle M, w \rangle \in A_{\text{TM}} \iff \langle M', w \rangle \in L_{\text{write}}.
    \]
    即 $A_{\text{TM}} \leq_m L_{\text{write}}$。由于 $A_{\text{TM}}$ 不可判定，故 $L_{\text{write}}$ 也不可判定。
\end{solution}


\question P152, 5.36题。
证明存在$\{1\}^*$的子集是不可判定的.
\begin{solution}
	设 $\Sigma = \{1\}$，则 $\Sigma^* = \{ \varepsilon, 1, 11, 111, \dots \} = \{1^n \mid n \ge 0\}$。

    \begin{enumerate}
        \item \textbf{可判定语言的数量：}
        每个可判定语言对应一个总停机的图灵机（判定器）。由于图灵机的描述是有限字母表上的有限串，因此图灵机的总数是可数无穷的。从而，可判定语言的总数至多为可数无穷。
    
        \item \textbf{$\{1\}^*$ 的子集数量：}
        $\{1\}^*$ 与自然数集 $\mathbb{N}$ 一一对应（$1^n \leftrightarrow n$），因此它的子集与 $\mathbb{N}$ 的子集一一对应。而 $\mathbb{N}$ 的子集有 $2^{\aleph_0}$ 个，这是不可数无穷的（康托尔定理）。
    
        \item \textbf{存在不可判定的子集：}
        如果 $\{1\}^*$ 的每个子集都是可判定的，那么 $\{1\}^*$ 的子集总数将不超过可判定语言的总数，即可数无穷。但事实上 $\{1\}^*$ 的子集有不可数无穷多个，因此必然存在 $\{1\}^*$ 的子集是不可判定的。
    
        \item \textbf{构造性证明（可选）：}
        虽然存在性证明已足够，但也可以显式构造一个不可判定的 $\{1\}^*$ 的子集。例如，取一个已知的不可判定语言 $B \subseteq \{0,1\}^*$（如 $A_{\text{TM}}$ 的编码集），然后通过一一映射 $\{0,1\}^* \to \{1\}^*$ 将 $B$ 映射到 $\{1\}^*$ 的一个子集。具体地，定义双射 $\varphi: \{0,1\}^* \to \{1\}^*$ 如下：
        \[
        \varphi(x) = 1^{\mathrm{num}(x)},
        \]
        其中 $\mathrm{num}(x)$ 是将 $x$ 视为二进制数（去掉前导零）对应的自然数。然后定义
        \[
        A = \{ \varphi(x) \mid x \in B \}.
        \]
        由于 $B$ 不可判定，且 $\varphi$ 是可计算的双射，则 $A$ 也是不可判定的（因为若有 $A$ 的判定器，则可通过 $\varphi^{-1}$ 得到 $B$ 的判定器）。
    \end{enumerate}
\end{solution}

\question P170, 6.3题，证明：如果$A\leq_T B $且 $B\leq_T C$，则 $ A\leq_T C$
\begin{solution}
	由 $A \leq_T B$ 知，存在谕示图灵机 $M_1^B$ 判定 $A$。  
由 $B \leq_T C$ 知，存在谕示图灵机 $M_2^C$ 判定 $B$。

我们需要构造一台谕示图灵机 $M^C$ 判定 $A$。构造思路如下：

$M^C$ 模拟 $M_1^B$ 的运行。每当 $M_1^B$ 查询谕示 $B$ 关于某个字符串 $q$ 是否属于 $B$ 时，$M^C$ 就调用 $M_2^C$（作为子程序）来判定 $q \in B$ 与否。由于 $M_2^C$ 可以判定 $B$，因此 $M^C$ 可以模拟 $M_1^B$ 对谕示 $B$ 的所有查询。

\noindent \textbf{具体构造：}
\begin{enumerate}
    \item 在输入 $x$ 上启动 $M_1$（即 $M_1^B$ 的不带谕示的普通图灵机部分）。
    \item 每当 $M_1$ 进入查询状态，准备询问字符串 $q$ 是否属于 $B$ 时：
        \begin{enumerate}
            \item 暂停 $M_1$ 的模拟。
            \item 以 $q$ 作为输入，运行 $M_2^C$（即 $M_2^C$ 的谕示版本，其谕示为 $C$）。
            \item 由于 $M_2^C$ 可以判定 $B$，它会在有限步内停机并给出“是”或“否”的答案。
            \item 将这个答案作为 $M_1$ 的谕示回答，恢复对 $M_1$ 的模拟。
        \end{enumerate}
    \item 继续模拟直到 $M_1$ 停机并输出结果。由于 $M_1^B$ 判定 $A$，这个输出就是 $x \in A$ 的正确判定。
\end{enumerate}

\noindent \textbf{正确性分析：}
\begin{itemize}
    \item 在 $M^C$ 的模拟中，每次对 $B$ 的查询都被 $M_2^C$ 正确回答（因为 $M_2^C$ 判定 $B$）。
    \item 因此，$M^C$ 的行为完全等同于 $M_1^B$，从而能够正确判定 $A$。
    \item $M^C$ 在谕示 $C$ 下总是停机，因为 $M_1^B$ 在谕示 $B$ 下总是停机（由 $A \leq_T B$ 保证），且每次调用 $M_2^C$ 也总是停机（由 $B \leq_T C$ 保证）。
\end{itemize}

构造出了谕示图灵机 $M^C$ 判定 $A$，即 $A \leq_T C$。
\end{solution}

\question P170, 6.7题。证明：对于任意的两个语言 $A$ 和 $B$，存在语言 $J$，使得 $A\leq _T J$且  $B\leq _T J$。
\begin{solution}
设 $A, B \subseteq \Sigma^*$。定义语言 $J$ 如下：
\[
J = \{0x \mid x \in A\} \cup \{1x \mid x \in B\}.
\]
这里我们在每个串前添加一个前缀 $0$ 或 $1$ 以区分来源。假设 $0$ 和 $1$ 不是 $\Sigma$ 中的符号（否则需要更谨慎的编码，但原理相同）。

\noindent \textbf{证明 $A \leq_T J$：}  
构造谕示图灵机 $M_1^J$ 判定 $A$：
\begin{enumerate}
    \item 输入 $x$。
    \item 查询谕示 $J$：“$0x \in J$ 吗？”
    \item 如果谕示回答“是”，则接受 $x$（即判定 $x \in A$）；否则拒绝 $x$。
\end{enumerate}
因为 $0x \in J \iff x \in A$，所以 $M_1^J$ 正确判定 $A$。因此 $A \leq_T J$。

\noindent \textbf{证明 $B \leq_T J$：}  
构造谕示图灵机 $M_2^J$ 判定 $B$：
\begin{enumerate}
    \item 输入 $x$。
    \item 查询谕示 $J$：“$1x \in J$ 吗？”
    \item 如果谕示回答“是”，则接受 $x$；否则拒绝 $x$。
\end{enumerate}
因为 $1x \in J \iff x \in B$，所以 $M_2^J$ 正确判定 $B$。因此 $B \leq_T J$。
\end{solution}

\question P171, 6.11题。证明 $\overline{\text{EQ}}_\text{TM}$可由一个带 $A_\text{TM}$的图灵机识别。
\begin{solution}
	\subsection*{证明思路}
    $L(M_1) \neq L(M_2)$ 当且仅当存在某个字符串 $x$ 使得 $x$ 属于 $L(M_1)$ 但不属于 $L(M_2)$，或者属于 $L(M_2)$ 但不属于 $L(M_1)$。即：
    \[
    L(M_1) \neq L(M_2) \iff \exists x \big[ (x \in L(M_1) \land x \notin L(M_2)) \lor (x \notin L(M_1) \land x \in L(M_2)) \big].
    \]
    由于 $A_{\text{TM}}$ 可以判定“$M$ 是否接受 $x$”，我们可以利用 $A_{\text{TM}}$ 谕示来搜索这样的 $x$。
    
    \subsection*{构造识别机}
    构造谕示图灵机 $M^{A_{\text{TM}}}$ 如下：
    
    $M^{A_{\text{TM}}}$ 对输入 $\langle M_1, M_2 \rangle$：
    \begin{enumerate}
        \item 枚举所有字符串 $x \in \Sigma^*$（按标准顺序，如短lex序）。
        \item 对每个枚举的 $x$，执行以下步骤：
        \begin{enumerate}
            \item 询问谕示 $A_{\text{TM}}$：$\langle M_1, x \rangle \in A_{\text{TM}}$ 是否成立？（即 $M_1$ 是否接受 $x$？）
            \item 询问谕示 $A_{\text{TM}}$：$\langle M_2, x \rangle \in A_{\text{TM}}$ 是否成立？（即 $M_2$ 是否接受 $x$？）
            \item 如果两个回答不同（即一个“是”一个“否”），则立即接受 $\langle M_1, M_2 \rangle$（因为找到了区分两个语言的串）。
        \end{enumerate}
        \item 如果对所有 $x$ 都没有找到区分串，则永不接受（即循环）。
    \end{enumerate}
    
    \subsection*{正确性分析}
    \begin{itemize}
        \item 如果 $\langle M_1, M_2 \rangle \in \overline{\text{EQ}}_{\text{TM}}$，即 $L(M_1) \neq L(M_2)$，则必然存在某个 $x$ 使得 $x$ 在 $L(M_1)$ 和 $L(M_2)$ 的成员关系不同。由于枚举了所有串，$M^{A_{\text{TM}}}$ 最终会枚举到这个 $x$，并发现两个谕示回答不同，从而接受输入。
        \item 如果 $\langle M_1, M_2 \rangle \notin \overline{\text{EQ}}_{\text{TM}}$，即 $L(M_1) = L(M_2)$，则对所有 $x$，两个谕示的回答总是相同，因此 $M^{A_{\text{TM}}}$ 永远不会接受，即它在输入上永不停机。
    \end{itemize}
    因此，$M^{A_{\text{TM}}}$ 识别了 $\overline{\text{EQ}}_{\text{TM}}$。
\end{solution}


 
\end{questions}
\end{document}

%%% Local Variables: 
%%% mode: latex
%%% TeX-master: t
%%% End: 
